% !TeX root = ../thuthesis-example.tex

\chapter{音系自进化生成与校准模型}
为实现一个能够通过用户互动进行自我优化的动态发音学习系统，本项目设计并实现了“音系自进化生成与校准模型”。该模型的核心思想是构建一个“\textbf{数据生成-质量过滤-模型微调}”的自动化闭环，旨在将静态的知识库，转化为一个具备持续学习与校准能力的动态系统。

\section{技术渊源与框架概览}
本模型的架构并非单一技术的简单应用，而是对近年来检索增强生成（RAG）领域多个前沿思想的有机融合与深度改造。其“技术基因”主要源于以下几个层面：
\begin{itemize}
    \item \textbf{基础范式}: 整体架构遵循了开创性的\textbf{RAG}范式，即一个结合了参数化内存（预训练语言模型）与非参数化内存（外部知识库）的系统。
    \item \textbf{检索策略}: 精密的检索模块，其两阶段设计思想源于\textbf{TeleOracle} 在处理专业领域（电信）知识时所采用的高级检索策略。
    \item \textbf{自进化机制}: 模型的自我迭代与数据增强循环，其核心的自动化数据生成与质量过滤流程，是对\textbf{CoRAG}方法论的创新性改编。
    \item \textbf{微调哲学与技术}: 模型的持续学习与优化过程，在训练理念上与\textbf{RAFT} (Retrieval-Augmented Fine-Tuning)的哲学一脉相承，并在具体实现上采用了高效的\textbf{LoRA} (Low-Rank Adaptation)技术。
\end{itemize}

本章将详细阐述，这些技术思想如何被整合并改造，以服务于本项目独特的多模态音系学习场景，从而构建出一个完整的自进化解决方案。
\section{音系知识驱动的跨模态检索与生成}
此模块是模型与用户进行交互的基础。在本模型中，非参数化内存特指“\textbf{多模态方言音系知识图谱}”与“\textbf{现代汉语方言音档}”，本模块架构图如\ref{fig:search}所示。

\begin{figure}[H]
    \centering
    \includegraphics[width=0.9\textwidth]{tele.pdf}
    \caption{检索策略架构图}
    \label{fig:search}
\end{figure}

\subsection{跨模态检索机制}
当用户输入一段发音练习录音 $x_u$ 时，从庞大知识库中精准地检索出最相关的“标准示范”是后续所有分析的前提。本项目设计的两阶段检索流程，是为了应对专业领域知识检索所提出的精细化策略。

\subsubsection{混合检索器}
此阶段的目标是快速、广泛地召回一个相关的候选集。专业领域既包含必须精确匹配的术语，也包含需要理解语义的复杂概念，因此必须采用混合检索。方言音系领域也具有此双重特性：词汇本身需要精确匹配，而发音的优劣则是一个声学上的相似度问题。

\paragraph{关键词检索实现：BM25}
为精确匹配词条，本项目采用经典的稀疏向量检索算法BM25。系统首先通过轻量级ASR模块获得用户发音的文本转写，并以此为查询，在“现代汉语方言音档”的文本描述中进行关键词匹配。其核心公式为：
\begin{equation}
\text{BM25\_score}(q, D) = \sum_{i=1}^{n} \text{IDF}(q_i) \cdot \frac{f(q_i, D) \cdot (k_1 + 1)}{f(q_i, D) + k_1 \cdot (1 - b + b \cdot \frac{|D|}{\text{avgdl}})}
\end{equation}

在本项目的应用中，$q_i$是用户发音转写文本中的词元，$D$是音档中的条目描述。例如，对于一个罕见的方言词汇，其$\text{IDF}(q_i)$值会很高，从而在检索中获得更高权重。参数$k_1$和$b$则通过在一个留存的方言语料验证集上进行网格搜索来优化，以达到最佳召回效果。

\paragraph{声学语义检索实现：Wav2Vec2与Faiss}
为捕捉发音本身的声学相似性，本项目采用了自监督学习在语音表征领域的杰出成果——Wav2Vec2 选用它的核心原因在于其自监督特性，使得本项目可以在标注数据稀缺的方言上，依然能获得高质量的声学表征。具体流程如下：用户的原始音频波形 $x_u$ 被输入到预训练的Wav2Vec2模型中，输出一系列帧级别的上下文声学表征。本项目对此序列施加一个均值池化层（mean-pooling），从而得到一个固定维度的、代表整个发音的稠密向量（dense vector）。

随后，该向量将作为查询，在利用Faiss (Facebook AI Similarity Search)构建的高效向量索引中进行最大内积搜索（MIPS）。本项目已预先将“多模态方言音系知识库”中所有标准发音的声学向量提取并构建成Faiss索引（根据规模可选用IndexFlatL2或IndexIVFPQ等索引类型），这使得本项目能在毫秒级时间内，从数以万计的标准音中，召回声学上最相似的top-N个候选。

\subsubsection{第二阶段：跨模态重排器（Cross-Modal Reranker）}
第一阶段混合检索的结果融合后，送入一个跨编码器（Cross-Encoder）架构的重排器进行精排。‘TeleOracle’的实践证明，在专业领域，牺牲一定的效率换取更高的精度是必要的。与在第一阶段中独立编码查询和文档的双编码器 (Bi-Encoder) 不同，跨编码器将用户发音$x_u$的特征与每一个候选标准示范$z_i$的特征进行拼接，作为一个整体对（pair）送入一个深度模型（如一个小型BERT）中进行联合编码。这使其能够捕捉到双编码器无法感知的、更深层次的交互信息。在本项目的场景中，这种交互信息至关重要，例如，它能直接建模用户发音在第二共振峰（F2）上的特定偏移，与某个候选元音的标准F2频率之间的差异。重排器最终为每个候选对 $(x_u, z_i)$ 输出一个范围在`[0, 1]`的精细相关性评分，系统据此选取唯一的、评分最高的标准示范$z^*$，跨编码器与双编码器架构对比如所示。
\begin{figure}[H]
    \centering
    \includegraphics[scale=0.3]{bicross.pdf}
    \caption{跨编码器与双编码器架构对比图}
    \label{fig:compare}
\end{figure}

\subsection{个性化辅导生成}
在获得最优标准示范$z^*$后，一个参数化的序列到序列模型 $\text{Generator}_{\theta}$ 开始工作，生成个性化反馈$a_d$。
输入给生成器的，是一个结构化的、融合了多模态信息的表示，例如：\texttt{[USER\_AUDIO\_VEC] [USER\_ASR\_TEXT] <SEP> [STANDARD\_AUDIO\_VEC] [STANDARD\_IPA\_TEXT]}。
\begin{equation}
    a_d = \text{Generator}_{\theta} (x_u, z^*)
\end{equation}
\section{自进化式数据增强与微调}
这是本模型实现“自进化”的核心机制，旨在将每一次用户互动都转化为一次潜在的模型学习机会。

\subsection{自进化核心流程：CoRAG与RAFT思想的融合}
本模块的整体流程，是对CoRAG自动化数据生成思想的改编，以及对RAFT监督微调思想的实践。CoRAG的核心是生成并筛选出最优的“检索链”来增强数据集，而本项目的模型则简化了这个过程：本项目生成的不是一个链，而是一个高质量的“诊断-示范-反馈”三元组 $(x_u, z^*, a_d)$。本项目使用“拒绝采样”思想来过滤这些三元组，然后采用RAFT的理念，用这些高质量、带“标准答案”（即$z^*$和$a_d$）的数据对模型进行微调。RAFT概览如\ref{fig:RAFToverview}所示。

\begin{figure}[H]
    \centering
    \includegraphics[width=0.9\textwidth]{RAFT概览.pdf}
    \caption{RAFT概览}
    \label{fig:RAFToverview}
\end{figure}

\subsection{高质量训练样本的生成与过滤}
每一次成功的交互都会构成一个潜在的训练样本三元组 $(x_u, z^*, a_d)$。为确保只有最优质的数据被用于训练，本项目引入了基于拒绝采样 (Rejection Sampling)的质量过滤机制。本项目设计了一个多维度的“\textbf{教师模型}” (Teacher Model) 来对每个生成的诊断反馈 $a_d$ 进行打分。该教师模型由两部分组成：
\begin{itemize}
    \item \textbf{逻辑教师 (Logical Teacher)}: 一个强大的、参数固定的LLM，负责评估反馈文本$a_d$的语言流畅度、逻辑清晰度以及解释的合理性。它会给出一个逻辑分数 $S_{\text{logic}} = P_{\text{LLM}}(a_d | x_u, z^*)$。
    \item \textbf{声学教师 (Acoustic Teacher)}: 一个专门训练的、高精度的声学诊断分类器。它直接分析用户音频$x_u$和标准音频$z^*$，并输出一个结构化的诊断结果（例如，`{error\_type: 'tone', detail: 'T3\_sandhi'}`）。然后，它会评估模型生成的文本反馈$a_d$是否与这个声学诊断结果一致，从而给出一个声学准确度分数$S_{\text{acoustic}}$。
\end{itemize}

最终的质量分是两者的加权和：$\text{score} = w_{\text{logic}}S_{\text{logic}} + w_{\text{acoustic}}S_{\text{acoustic}}$。只有当分数超过阈值$\tau$时，该样本才被“接受”进入微调数据集 $\mathcal{D}_{\text{finetune}}$。

\subsection{监督微调：RAFT哲学的实践与LoRA实现}
\subsubsection{训练哲学：融合RAFT思想}
当微调数据集$\mathcal{D}_{\text{finetune}}$积累到一定规模后，系统会自动启动SFT流程。此过程完全体现了\textbf{RAFT}的核心哲学：训练模型学会“利用参考资料进行思考”。如\ref{fig:RAFTfine}在本项目的场景中，每一个训练样本 $(x_u, z^*, a_d)$ 都可以被看作一个完美的RAFT式实例：
\begin{itemize}
    \item \textbf{问题 (Question)}: 用户的发音 $x_u$ (一个需要被诊断的问题)。
    \item \textbf{参考资料 (Oracle Document)}: 检索到的标准示范 $z^*$ (解决问题所必需的权威依据)。
    \item \textbf{CoT式答案 (Chain-of-Thought Answer)}: 模型生成的诊断反馈 $a_d$ (一段解释了“如何从问题到答案”的、带有推理过程的文本)。
\end{itemize}
通过在这种高质量、带上下文的“纠错对”上进行训练，模型学习到的不仅仅是语言模型，更是一种“音系诊断”的专业技能。

\begin{figure}[H]
    \centering
    \includegraphics[width=0.9\textwidth]{RAFT微调.pdf}
    \caption{RAFT微调}
    \label{fig:RAFTfine}
\end{figure}

\subsubsection{实现技术：高效的LoRA微调}
考虑到全量微调大型生成模型的计算成本极高，本项目采用了低秩自适应 (LoRA)技术。LoRA的假设是，在微调过程中，模型权重的改变量 $\Delta W$ 是低秩的。因此，它可以用两个更小的、低秩的矩阵 `B` 和 `A` 的乘积来近似，即 `$\Delta W \approx B \cdot A$`。最终，模型权重的更新公式变为：
\begin{equation}
    W_{\text{new}} = W_0 + \frac{\alpha}{r} B \cdot A
\end{equation}

其中，$W_0$是冻结的原始预训练权重，$B \in \mathbb{R}^{d \times r}$ 和 $A \in \mathbb{R}^{r \times k}$ 是可训练的低秩适配器矩阵，其秩$r$远小于原始维度$d, k$。在本项目的实现中，本项目通常会选择一个很小的秩（如$r=8$或$16$），并将LoRA适配器应用于生成器模型中Transformer架构的注意力层的查询（Query）和值（Value）矩阵上。这使本项目仅用不到1\%的额外参数，就能达到接近全量微调的效果。

SFT的损失函数旨在最大化模型在给定用户发音和标准示范条件下，生成高质量诊断反馈的对数似然：
\begin{equation}
    \mathcal{L}_{\text{SFT}}(\theta_{\text{LoRA}}) = - \sum_{(x_u, z^*, a_d) \in \mathcal{D}_{\text{finetune}}} \log P_{\theta_0 + \theta_{\text{LoRA}}}(a_d | x_u, z^*)
\end{equation}

通过此过程，模型$\text{Generator}_{\theta}$的LoRA适配器参数得以持续更新，使其能不断从真实用户互动中汲取养分，提升诊断的准确性和反馈的质量，最终形成一个完整的、自动化的“\textbf{交互-生成-过滤-学习}”的自进化闭环。
